\scanpage{5}
\mspage{4}
utile d'introduire d'autres variantes.

3) \deleted{\unsure{S}} Surface obtenue comme $\mathcal X_0$ en excluant les p.r. super\deleted{singulières}\inserted{critiques}, mais en introduisant, pour éviter les sommets supercritiques, des arêtes supplémentaires du graphe \inserted{$\Gamma$} des p.r., savoir les ``balayages au-dessous d'une p.r. supercritique'' \inserted{$D_i$}, et ceci aussi bien pour une p.r. voisine sous-stable, que stable. Cela fait donc des arcs de transition supplémentaires, qui \deleted{\unsure{remplacent}}\inserted{prennent} la place des arcs de spéc. supercritique de $2^\circ$, mais qui ne sont plus des arcs de spécialisation, à proprement parler. Et d'ailleurs l'opération correspondante doit être considérée comme étant \uline{de poids $2$}, et non $1$ comme les autres.
\figure{%
\begin{tikzpicture}[x=1cm,y=1cm,line cap=round]
\draw[yellow!65!black,line width=3pt] (-3,0)..controls(-1,.12)and(1,.2)..(3,.28);
\draw[yellow!65!black,line width=3pt] (-2.55,.75)--(-2.35,-.65) (-1.5,1.05)--(-1.25,-.65) (.0,1.08)--(.03,-.6) (1.2,1.02)--(1.2,-.62) (2.5,.96)..controls(2.55,.4)and(2.5,-.3)..(2.43,-.58);
\end{tikzpicture}%
}
On trouve, dans la surface $\mathcal X_2$ obtenue, \deleted{et} (décomposée \inserted{cellulairement} par le graphe dual de $\Gamma$) des sommets qui à nouveau sont tous d'ordre $4$, ils \inserted{(ceux qui sont proches \deleted{\unsure{à}} $D_i$)} correspondent aux \deleted{\unsure{triplets de deux arêtes}} \inserted{paires $(s,a)$ où $s$ \unsure{sommet de} $D_i$, $a$ arête sur $D_i$ incidents} de $D_i$ \deleted{\unsure{d'une arête}\illegible{3}\unsure{adjacentes}}, i.e. aux arêtes de la subdivision barycentrique \deleted{de $\unsure{X}$} de $D_i$, il y en a $2(n-1)$ dans le cas stable.
Si je me rappelle bien, $D_i$ est plongée dans $\mathcal X_3$, de faç que les sommets de la \deleted{\unsure{subdivision}}\inserted{décomposition} cell. de $D_i$ double de la subd. bar. de celle qu'on a au départ, \deleted{\unsure{tous}} correspondent à des sommets de $\mathcal X_3$.
\margin{4}
